% Owned proposed fragment B.130; central integration is handled by root.
\subsection*{An actual source relation and a retracted two-dimensional torus group}

\statement{B.130}{Proved}{The actual six-generator one-neck source map is noninjective in both categories, and its image contains an explicitly retracted copy of $\mathbb Z^2$}
\label{stmt:B130}
Let $\mathcal G=\mathcal G^0_{\rm ann}$ be the actual ordinary
one-neck group of B.114, B.119 and B.123. Use its certified common
marking $K=F(a,B,d,e,f)$ and selected plus gap cone $\tau$.
Let
\[
 (s_1,s_2,s_3,s_4,v,m)
\]
be the actual ordered six source images, with $s_i$ the four
outer contact words and $m=k_-$ the inner good-region source.
Write $\delta=\operatorname{Inn}(\tau)$ and put
\begin{equation}\label{eq:B130-U}
 U=s_1s_4^{-1}s_1^{-1}.
\end{equation}
Then the following are exact identities in the actual group:
\begin{equation}\label{eq:B130-actual-relation}
 U=a\tau a\tau,\qquad [m,U]=1.
\end{equation}
The source word
\begin{equation}\label{eq:B130-source-word}
 m s_1s_4^{-1}s_1^{-1}m^{-1}s_1s_4s_1^{-1}
\end{equation}
is nontrivial in the abstract free source group $F_6$ and in its
full profinite completion. Its image is trivial in $\mathcal G$
and in $\widehat{\mathcal G}$, respectively.

Moreover $\langle U,m\rangle$ is an actual retract isomorphic to
$\mathbb Z^2$ in $\mathcal G$ and in the actual source image $H$.
Its closure in $\widehat{\mathcal G}$ is a continuous retract
canonically isomorphic to $\widehat{\mathbb Z}^2$, and the
same retraction restricts to the actual closed source image
$\overline H$. These conclusions do not decide the equality
of that source image with the selected fixed subgroup.

There is also a specified nontrivial element in the selected
centralizer of the actual intermediate full completion
$\widehat M$ of B.120 which is killed by the final four-prefix
filling. Thus this filling is noninjective on that completed
intermediate centralizer; no embedding of an abstract
$\widehat F_6$ into $\widehat M$ is needed.

\begin{proof}
\emph{B.130.1. The actual common marking and three node relators.}
Let $\eta$ be the actual based mirror of B.123 and define
\begin{equation}\label{eq:B130-local}
 A=\eta(a),\qquad B_+=\eta(B),\qquad
 C=A^{-1}\tau A,\qquad b=CB_+C,\qquad V=AB_+A^{-1}.
\end{equation}
The physical prefix formulas on $O_1=aBa^{-1}$ and $O_2=a$
give, as literal free words,
\begin{equation}\label{eq:B130-mirror-identities}
 V=\eta(O_1)=aB^{-1}a^{-1},\qquad
 A=V^{-1}a^{-1}V,\qquad VA^{-1}V^{-1}=a.
\end{equation}
The plus regional presentation includes the three actual node
relations
\begin{equation}\label{eq:B130-three-nodes}
 r=[A,CAC]=1,\qquad [b,CbC]=1,\qquad [A,b]=1.
\end{equation}
There is no change of their based meridian convention. Conjugation
by $A$, inversion where indicated, and conjugation by $\tau$
give the literal identities
\begin{equation}\label{eq:B130-four-crosses}
\begin{aligned}
 ArA^{-1}&=[A,\delta(A)],\\
 \bigl(A[b,CbC]A^{-1}\bigr)^{-1}&=[V,\delta(V)],\\
 A[A,b]A^{-1}&=[A,\delta(V)],\\
 \delta\bigl(A[A,b]A^{-1}\bigr)^{-1}&=[V,\delta(A)].
\end{aligned}
\end{equation}
Consequently the subgroups $D=\langle V,A\rangle$ and
$\delta(D)$ commute elementwise. The map
\[
 N_\delta(u)=u\delta(u)
\]
is a homomorphism on $D$, since
$u\delta(u)v\delta(v)=uv\delta(u)\delta(v)$ for $u,v\in D$.
This uses only the four displayed cross commutations; no
freeness, normality or simultaneous-basis assertion about $D$
is needed.

\emph{B.130.2. The raw source words and the node correction.}
The actual raw table of B.119, with the entire outer gap
conjugation retained, reads
\begin{equation}\label{eq:B130-raw-sources}
\begin{aligned}
 s_1&=A(CbCb)A^{-1},\\
 s_4&=A\bigl(Cb[A,CAC]b^{-1}C\,ACAC\bigr)A^{-1}.
\end{aligned}
\end{equation}
Put $q=ACb$. Free reduction using only the cone squares gives
\begin{equation}\label{eq:B130-raw-correction}
 s_1=N_\delta(V),\qquad
 s_4=(qrq^{-1})N_\delta(A).
\end{equation}
The first relation in \eqref{eq:B130-three-nodes} therefore gives
$s_4=N_\delta(A)$ in the actual group. Multiplicativity on $D$
and \eqref{eq:B130-mirror-identities} now give
\[
 U=N_\delta(V)N_\delta(A)^{-1}N_\delta(V)^{-1}
   =N_\delta(VA^{-1}V^{-1})=N_\delta(a).
\]
Finally the actual inner generic relation is $[m,a]=1$, and
the actual cross relation is $[\tau,m]=1$. Thus $m$ commutes
with $N_\delta(a)=a\tau a\tau$, proving
\eqref{eq:B130-actual-relation}. The proof takes place in the
full defining presentation, so no unproved survival of a
regional vertex group is involved.

\emph{B.130.3. Nontriviality in the actual source domain.}
The word \eqref{eq:B130-source-word} is freely and cyclically
reduced of length eight in the marked source basis. It is
therefore nontrivial in $F_6$. One finite quotient also proves
its nontriviality in $\widehat F_6$ without any assertion about
the topology of its image. In
$\operatorname{UT}_3(\mathbb F_2)$ use coordinates
\[
 (u,v,z)(u',v',z')=(u+u',v+v',z+z'+uv').
\]
Map $m$ to $(1,0,0)$, $s_4$ to $(0,1,0)$ and the other four
source generators to the identity. The image of
\eqref{eq:B130-source-word} is $(0,0,1)$, the nonidentity
central element. The actual group relation proves that the
same nontrivial element of $\widehat F_6$ is killed by the
completed source map.

\emph{B.130.4. Two actual characters and the retraction.}
By B.124 the actual outer subgroup
\[
 H_{\rm out}=F(s_1,s_2,s_3,s_4,v)
\]
is a retract of $\mathcal G$, and that retraction sends $m$ to
$1$. Compose it with minus the $s_4$ exponent to obtain
$\chi_1:\mathcal G\to\mathbb Z$. The actual base projection
of B.119 and B.123, followed by the $m$ exponent, gives
$\chi_2:\mathcal G\to F(m,v)\to\mathbb Z$. Then
\[
 (\chi_1(U),\chi_1(m))=(1,0),\qquad
 (\chi_2(U),\chi_2(m))=(0,1).
\]
Thus $r=(\chi_1,\chi_2):\mathcal G\to\mathbb Z^2$ has the
section
\[
 i:\mathbb Z^2\longrightarrow\mathcal G,
 \qquad i(n,\ell)=U^n m^\ell.
\]
It is a homomorphism by \eqref{eq:B130-actual-relation}, and
$ri=\mathrm{id}$. This proves the asserted embedding and
retraction. Its image lies in $H$, so restriction proves the
same assertion for that actual source image.

On full completions $\widehat r\widehat i=\mathrm{id}$.
Therefore $\widehat i$ is injective and its image is exactly
the closure of $\langle U,m\rangle$, since the discrete subgroup
is dense in $\widehat{\mathbb Z}^2$. This gives its full
induced topology and the continuous retraction. Restricting to
$\overline H$ gives the final claimed retraction. No equality
between the abstract completion of all of $H$ and
$\overline H$ has been used.

\emph{B.130.5. A nonzero completed intermediate-centralizer element.}
Use the actual intermediate proper amalgam of B.120,
\[
 M=R*_{H_\partial}W_\partial,\qquad
 \widehat M=\widehat R
       \amalg_{\widehat H_\partial}\widehat W_\partial.
\]
Write \(c\) for its selected regional cone and
\(\ell_1,\ldots,\ell_5\) for the five actual selected-gap
regional source words. These generate the actual free-five
source subgroup with its full induced topology by B.112.
The symbols \(\ell_i\) distinguish those intermediate words
from the final \(s_i\). Put
\[
 \ell=\ell_1\ell_4^{-1}\ell_1^{-1},\qquad
 D=jcj,\qquad z=[k,j\ell j].
\]
Both factors centralize \(D\) in \(M\): the first does by the
actual corner cross relation, and the second because
\(\ell\) centralizes \(c\) in \(R\). Under the actual final map
\[
 M\longrightarrow
 M/\langle\!\langle(jP_i)^2:1\leq i\leq4\rangle\!\rangle
\]
the marked words \(j\ell_i j\) map to \(s_i\) for \(i\leq4\),
\(j\ell_5j\) maps to \(v\), and \(k\) maps to \(m\).
Therefore \(z\) maps to \([m,U]=1\).

We verify its nontriviality even in \(\widehat M\) by a finite
separation argument. Conjugate the proper amalgam by \(j\),
which belongs to its corner vertex:
\[
 \widehat M=\widehat R^{\,j}
       \amalg_{\widehat H_\partial^{\,j}}\widehat W_\partial .
\]
First \(j\ell j\notin\widehat H_\partial^{\,j}\).
Otherwise \(\ell\in\widehat H_\partial\) would centralize \(c\).
The exact edge group is the free profinite product
\[
 \widehat H_\partial=C_2(c)\amalg\widehat{\mathbb Z}^{\,2},
\]
whose \(c\)-centralizer is \(\langle c\rangle\) by the
factor-centralizer theorem used in B.118--B.120.
The source word \(\ell\) has original orientation zero, whereas
\(c\) has orientation one. It would follow that \(\ell=1\),
contradicting its nontrivial free-five word and full regional
source topology.

Second \(k\notin\widehat H_\partial^{\,j}\). Indeed
\(k\in\widehat H_\partial\), and the actual completed corner
face of B.126 is
\[
 \widehat H_\partial\cap j\widehat H_\partial j
       =\overline{\langle P\rangle},\qquad P=cT.
\]
The finite character on the edge which sends \(k\) to the
nontrivial element of \(C_2\) and sends \(c,T\) to \(1\)
kills this entire procyclic subgroup, but does not kill \(k\).

Each of these exclusions can be retained in a finite quotient
of its compact vertex group, because the embedded edge is
closed. B.120's cofinal compatible-kernel construction,
transported by \(j\), refines the two prescribed finite
targets. It gives finite vertex images \(F_R,F_W\), with the
same edge image \(E\) embedded in both, in which the images of
\(j\ell j\) and \(k\) still lie outside \(E\).
Their commutator is a reduced word of length four in the
ordinary amalgam \(F_R*_E F_W\).

For completeness, that finite-vertex amalgam is residually
finite. Align sufficiently many regular permutation actions
of the two finite vertices to agree on \(E\). The resulting
finite permutation representation is injective on both
vertices. Its kernel acts freely on the Bass--Serre tree and
is a free group; it has finite index and is finitely generated.
A nontrivial kernel element is separated by a finite quotient
of that free group, whose finite-index core in the whole
amalgam supplies a finite quotient of the amalgam. Elements
outside the kernel are already detected by the permutation
representation. Thus a finite target detects our reduced
commutator. Composing with the actual map from \(M\) shows
that \(z\ne1\) in \(\widehat M\).

We have proved a nonzero element of
\(C_{\widehat M}(D)\) is killed by the actual final filling.
This argument uses compatible finite targets and their
ordinary normal forms, not an unsupported normal-form
assertion inside an arbitrary profinite amalgam.

\emph{B.130.6. Exact evidence and scope.}
The input-bound certificate checks all raw identities above by
free cancellation and cone squares, retains the exact conjugates
of all three node relators, checks both signed actual $m$ actions,
and evaluates the single order-eight source detector. The
multiplicativity step is also checked in the universal group
$F(V,A)\times F(\delta V,\delta A)$ presented by the four cross
commutators. This is a complete algebraic implication, not a
finite quotient test for equality in the actual group.

The earlier one-sided source detector and its undecided full
free-word residue remain preserved. The argument here adds the
actual node relations that were missing from that detector.
It proves noninjectivity of this six-generator source map;
it does not exhibit a fixed element outside its source image
or resolve either open equality in B.123.
\end{proof}

\begingroup\raggedright\footnotesize
Exact certificate and genuine outputs:
\path{work_CE/code/check_actual_source_commutator26.py},
\path{work_CE/code/check_actual_source_commutator26.out},
\path{work_CE/code/check_actual_source_commutator26.err};
\path{work_CE/next_symmetric_nodes/actual_source_commutator26.json}.
The full research proof is
\path{work_CE/actual_source_commutator26.txt}.
The complete independent finite-separation proof is
\path{work_A/intermediate_completed_source_relator_nontrivial26.md}.
\par\endgroup
